\begin{footnotesize}

% ========================================
% Table 1: Parameters from Literature
% ========================================
\begin{longtable}{llcc}
\caption{Model Calibration: Parameters from Literature} \label{tab:calibration_literature} \\
\toprule
\textbf{Parameter} & \textbf{Description} & \textbf{Value} & \textbf{Source} \\
\midrule
\endfirsthead
\multicolumn{4}{c}{\tablename\ \thetable\ -- \textit{Continued from previous page}} \\
\toprule
\textbf{Parameter} & \textbf{Description} & \textbf{Value} & \textbf{Source} \\
\midrule
\endhead
\midrule
\multicolumn{4}{r}{\textit{Continued on next page}} \\
\endfoot
\bottomrule
\endlastfoot
\multicolumn{4}{l}{\textit{Production Technology}} \\
$\epsilon$ & Elast. subst., varieties & 10 & \cite{ferrante2023inflationary} \\
\midrule
\multicolumn{4}{l}{\textit{Price Setting}} \\
$\kappa_i$ & Price adj. cost, sector $i$ & Table \ref{tab:sectoral_params} & \cite{pasten2021sectoral} \\
\midrule
\multicolumn{4}{l}{\textit{Household}} \\
$\gamma$ & Inv. intertemporal elast. & 2 & \cite{ferrante2023inflationary} \\
$\psi$ & Inv. Frisch elasticity & 1 & \cite{ferrante2023inflationary} \\
$\chi$ & Labor disutility weight & 1 & \cite{ferrante2023inflationary} \\
\midrule
\multicolumn{4}{l}{\textit{Monetary Policy}} \\
$\phi_\pi$ & Taylor rule coef. (inflation) & 2.5 & Literature \\
$\rho_i$ & Interest rate smoothing & 0.74 & Literature \\
\end{longtable}

\vspace{1em}

% ========================================
% Table 2: Parameters from Data
% ========================================
\begin{longtable}{llcc}
\caption{Model Calibration: Parameters from Data} \label{tab:calibration_data} \\
\toprule
\textbf{Parameter} & \textbf{Description} & \textbf{Value} & \textbf{Source} \\
\midrule
\endfirsthead
\multicolumn{4}{c}{\tablename\ \thetable\ -- \textit{Continued from previous page}} \\
\toprule
\textbf{Parameter} & \textbf{Description} & \textbf{Value} & \textbf{Source} \\
\midrule
\endhead
\midrule
\multicolumn{4}{r}{\textit{Continued on next page}} \\
\endfoot
\bottomrule
\endlastfoot
\multicolumn{4}{l}{\textit{Production Technology}} \\
$\alpha_{mi}$ & Material share, sector $i$ & Table \ref{tab:sectoral_params} & I-O Tables (Subsection \ref{material_share_IO}) \\
$\alpha_{vi}$ & Import share, sector $i$ & Table \ref{tab:sectoral_params} & I-O Tables (Subsection \ref{material_share_IO}) \\
$\Gamma_{i,j}$ & Input-output matrix & Table \ref{tab:io_matrix} & Chilean I-O Tables (2021) \\
\midrule
\multicolumn{4}{l}{\textit{Consumption}} \\
$\Omega_G$ & Goods weight in consumption & 0.57 & National Accounts (Subsection \ref{cons_share}) \\
$\Omega_S$ & Services weight in consumption & 0.43 & National Accounts (Subsection \ref{cons_share}) \\
$\omega_i^g$ & Goods sector $i$ weight & Table \ref{tab:consumption_shares} & National Accounts (Subsection \ref{cons_share}) \\
$\omega_i^s$ & Services sector $i$ weight & Table \ref{tab:consumption_shares} & National Accounts (Subsection \ref{cons_share}) \\
$\varrho_i$ & Home bias, sector $i$ & Table \ref{tab:sectoral_params} & Trade Data (Subsection \ref{home_bias}) \\
\midrule
\multicolumn{4}{l}{\textit{Small Open Economy}} \\
$P_t^*$ & Foreign price level (SS) & 1 & Normalized \\
$Y_t^*$ & Foreign output (SS) & 1 & Normalized \\
$\Pi_t^*$ & Foreign inflation (SS, quarterly) & 1.01 & International Data \\
$R^w$ & World interest rate (SS, quarterly) & 1.045 & International Data \\
\end{longtable}

\vspace{1em}

% ========================================
% Table 3: Calibrated Parameters
% ========================================
\begin{longtable}{llcc}
\caption{Model Calibration: Calibrated and Estimated Parameters} \label{tab:calibration_calibrated} \\
\toprule
\textbf{Parameter} & \textbf{Description} & \textbf{Value} & \textbf{Method} \\
\midrule
\endfirsthead
\multicolumn{4}{c}{\tablename\ \thetable\ -- \textit{Continued from previous page}} \\
\toprule
\textbf{Parameter} & \textbf{Description} & \textbf{Value} & \textbf{Method} \\
\midrule
\endhead
\midrule
\multicolumn{4}{r}{\textit{Continued on next page}} \\
\endfoot
\bottomrule
\endlastfoot
\multicolumn{4}{l}{\textit{Production Technology}} \\
$\epsilon_Y$ & Elast. subst., production inputs & 0.8 & Estimated \\
$\epsilon_{mi}$ & Elast. subst., materials & 0.1 & Calibrated \\
\midrule
\multicolumn{4}{l}{\textit{Price Setting}} \\
$\kappa_v$ & Price adj. cost, imports & $10^{13}$ & Calibrated (flexible prices) \\
\midrule
\multicolumn{4}{l}{\textit{Household}} \\
$\beta$ & Discount factor (quarterly) & 0.986 & Calibrated to match interest rate \\
\midrule
\multicolumn{4}{l}{\textit{Consumption}} \\
$\sigma_H$ & Elast. subst., home-foreign goods & 0.999 & Calibrated \\
\midrule
\multicolumn{4}{l}{\textit{Labor}} \\
$\Psi_l$ & Aggregate labor adj. cost & 10 & Calibrated \\
$\psi_l$ & Sectoral labor adj. cost & 0 & Calibrated (flexible) \\
\midrule
\multicolumn{4}{l}{\textit{Small Open Economy}} \\
$\eta^*$ & Export demand elasticity & 1 & Calibrated \\
$\chi_{i}^X$ & Export share, sector $i$ & $1/12$ & Calibrated (uniform) \\
$\omega^x$ & Foreign demand scale & 1 & Calibrated \\
\midrule
\multicolumn{4}{l}{\textit{Financial}} \\
$\xi_b$ & Debt elasticity (risk premium) & 0.001 & Calibrated \\
$\bar{b}$ & Steady-state debt-to-GDP & $-0.57$ & Endogenous (from TB target) \\
\midrule
\multicolumn{4}{l}{\textit{Shocks: Preferences}} \\
$\rho_{om,1}$ & Preference shock, AR(1) coef. & 0.1 & Calibrated \\
$\sigma_{om}$ & Preference shock, std. dev. & 0.0001 & Calibrated \\
\midrule
\multicolumn{4}{l}{\textit{Shocks: Technology}} \\
$\rho_{tfp,1}$ & TFP shock, AR(1) coef. & 0.5 & Calibrated \\
$\sigma_{tfp}$ & TFP shock, std. dev. & 0.001 & Calibrated \\
\midrule
\multicolumn{4}{l}{\textit{Shocks: Monetary Policy}} \\
$\sigma_i$ & Monetary shock, std. dev. & 0.001 & Calibrated \\
\end{longtable}

\end{footnotesize}

% ========================================
% Descriptive Text
% ========================================

\paragraph{Literature-based parameters (Table \ref{tab:calibration_literature})} 
Standard macroeconomic parameters follow \citet{ferrante2023inflationary}: the elasticity of substitution between varieties ($\epsilon=10$), risk aversion coefficient ($\gamma=2$), inverse Frisch elasticity ($\psi=1$), and labor disutility weight ($\chi=1$). Sectoral price adjustment costs ($\kappa_i$) are based on price flexibility estimates from \citet{pasten2021sectoral}. Monetary policy rule coefficients ($\phi_\pi=2.5$, $\rho_i=0.74$) are consistent with estimates for small open economies in the literature.

\paragraph{Data-based parameters (Table \ref{tab:calibration_data})} 
Production structure parameters are derived from Chilean Input-Output tables (2021, 12-sector aggregation). The input-output matrix $\Gamma_{i,j}$ and sectoral material shares ($\alpha_{mi}$, $\alpha_{vi}$) come directly from these tables (Subsection \ref{material_share_IO}). Consumption parameters ($\Omega_G$, $\Omega_S$, $\omega_i^g$, $\omega_i^s$) are constructed from National Accounts data on final consumption expenditure by product (Subsection \ref{cons_share}). Home bias parameters ($\varrho_i$) are calculated from trade statistics showing the domestic share of consumption by sector (Subsection \ref{home_bias}). Foreign economy steady-state values reflect international data: quarterly inflation of 1\% ($\Pi^*=1.01$) and a real interest rate implying 4.5\% annual rate ($R^w=1.045$).

\paragraph{Calibrated parameters (Table \ref{tab:calibration_calibrated})} 
Several structural parameters are calibrated to match key features of the Chilean economy. The discount factor ($\beta=0.986$) is set to match the steady-state interest rate. The elasticity of substitution in production ($\epsilon_Y=0.8$) is estimated to match sectoral output volatilities. Import price adjustment costs are set extremely high ($\kappa_v=10^{13}$) to approximate flexible import prices. Labor adjustment costs are calibrated with aggregate frictions ($\Psi_l=10$) but no sectoral frictions ($\psi_l=0$). Trade elasticities ($\eta^*=1$, $\epsilon_{mi}=0.1$) and export shares ($\chi_i^X=1/12$) are calibrated jointly with the foreign demand scale ($\omega^x=1$) to match Chile's trade openness and sectoral export patterns. The debt elasticity parameter ($\xi_b=0.001$) determines the sensitivity of the risk premium to foreign debt deviations. Steady-state debt ($\bar{b}=-0.57$) is endogenously determined by the trade balance target (2\% of GDP deficit). Shock processes (preference, TFP, and monetary policy) are calibrated to replicate observed business cycle dynamics, with persistence and volatility parameters chosen to match Chilean aggregate and sectoral data moments.
